\documentclass[10pt,a4paper]{amsart}
\usepackage[margin=19mm]{geometry}
\usepackage{amsmath,amssymb,mathtools}
\usepackage[scr=rsfs,cal=euler]{mathalfa}
\usepackage{xfrac}
\usepackage[unicode,hidelinks,bookmarks=false]{hyperref}
\newcommand{\ie}{i.e.\ }
\newcommand{\cf}{cf.\ }
\newcommand{\resp}{respectively\ }
\newcommand{\Subsection}{\subsection}
\providecommand{\nobreakdash}{\nobreak\hbox{-}}
\setlength{\emergencystretch}{2em}
\multlinegap0pt
\usepackage{amssymb}
\newtheorem{thm}{Theorem}
\title{Image Recognition via Vaisman--Neifeld's Geometry}
\author{N. C. Combe, H. K. Nencka}
\date{Source: arXiv:2607.26749v1 (CC BY 4.0)}
\begin{document}
\maketitle

\noindent\emph{Source and licence.} Image Recognition via Vaisman--Neifeld's Geometry. Version arXiv:2607.26749v1 (CC BY 4.0), distributed by arXiv under the Creative Commons Attribution 4.0 International licence, http://creativecommons.org/licenses/by/4.0/, as recorded in the arXiv licence metadata for this exact version and shown on the abstract page https://arxiv.org/abs/2607.26749v1. Full licence notice: \texttt{licenses/arxiv-2607.26749-CC-BY-4.0.txt}.\par\smallskip

\noindent\textit{Editorial scope.} Counted unit: the article's thm P:2 (source0.tex lines 253--279). The article's complete original proof of it is delivered with it (source0.tex lines 282--368, one \texttt{proof} environment of 87 lines). No mathematics is written, repaired or re-derived: the statement and the proof are the article's own lines, in the article's own notation, and no retained line is substituted or re-flowed.  Retained uncounted as the source-local context the proof needs: lines 139--149.  Nothing was omitted from any retained span, so the package carries no omission record.  No retained line contains a citation, so the wrapper carries no bibliography and no bibliography entry.  No retained line contains a figure environment, an image-inclusion command or drawing code, so no image is delivered and the package declares no asset dependency.  Editorial formatting only, in addition: an AMS \texttt{amsart} preamble that restates the article's own declarations and macros used by the retained lines, namely the environments \texttt{thm} on separate counters, together with the standard packages those lines need.  The retained environments print in this document as Theorem 1 in order of appearance, whereas the article prints the same items with the numbering of its own full document.  The complete native source of the article as distributed in the arXiv e-print for this exact version is bundled unchanged as \texttt{sources/206292/source0.tex}.
\begin{thm}\label{P:1} Let \( O \subset \mathbb{CP}^3 \) be a smooth, non-symmetric 3D object (real 3-dimensional submanifold), and let \( \Pi_1, \Pi_2: \mathbb{CP}^3 \dashrightarrow \mathbb{CP}^2 \) be two independent rational projections onto distinct complex projective planes. Assume:

\begin{enumerate}
\item  {\bf Non-degeneracy of projections}: The restrictions of \( \Pi_1 \) and \( \Pi_2 \) to \( O \) are immersive (the differentials \( d\Pi_i \) have full rank).

\item  {\bf Involution symmetry:} The projections satisfy \(\Pi_2 = \iota \circ \Pi_1 ,\) where \( \iota: \mathbb{CP}^2 \to \mathbb{CP}^2 \) is an anti-holomorphic involution (e.g., a polarity induced by the Fubini-Study metric).

\end{enumerate}

Each projection \( \Pi_i \) defines a holomorphic line bundle \( L_i \to \mathbb{CP}^2 \) equipped with a connection \( \nabla_i \) derived from the Fubini-Study metric. Then, the  involution \( \iota \) induces a duality \[ L_1 \leftrightarrow L_2^* ,\] making \( (\nabla_1, \nabla_2) \) a dual pair.
\end{thm}

\begin{thm}\label{P:2}
Let \( \mu_i: O \to \mathbb{C}^2 \) be the first moment maps of \( O \) for projections \( \Pi_i \), encoding the centroids of the projected data, and $\nabla_i$ Chern conections on bundles $E_i=\Pi_i^*T\mathbb{CP}^2$ with curvatures $F_{\nabla_i}$.

\smallskip 

Under these conditions:

\begin{itemize}
\item  {\bf Uniqueness}: The orientation of \( O \) (i.e., its position modulo projective transformations) is uniquely determined by the compatibility of \( \nabla_1 \) and \( \nabla_2 \) acting on \( \mu_1 \) and \( \mu_2 \).

\item {\bf Reconstruction}: The original object \( O \) can be reconstructed as the intersection of the parallel transports along \( \nabla_1 \) and \( \nabla_2 \), applied to the moment maps \( \mu_1 \) and \( \mu_2 \).
\end{itemize}
Explicitly, there exists a unique solution \( v \in T\mathbb{CP}^3 \), a {\bf direction vector} (modulo scaling) in the ambiant space, satisfying:

\[
\begin{cases}
\nabla_1 \mu_1 = v \cdot \omega_1, \\
\nabla_2 \mu_2 = v \cdot \omega_2, 
\end{cases}
\]

where \( \omega_i \) are {\bf connection 1-forms} encoding the involution duality \( \omega_1 = \iota^* \omega_2 \);
 $\nabla_i \mu_i$ is the covariant derivative of $\mu_i$, resulting in a $\mathbb{C}^2$-valued 1-form on $O$ \footnote{ The symbol $``\cdot"$ denotes the contraction of the vector $v$ with the $\mathbb{C}^2$-valued  1-form $\omega_i$. Concretely, 
$v \cdot \omega_i$ is the vector in $\mathbb{C}^2$ obtained by evaluating each $\omega_i^j$ on $v$: 

\[ v \cdot \omega_i=\bigg(\begin{smallmatrix} \omega_i^1(v)\\ \omega_i^2(v)\end{smallmatrix}\bigg)\]}. 
\end{thm}

\begin{proof}
\begin{enumerate}
    \item {\it Orientation and Duality of Connections}:
We prove that the involution $\iota$ forces the connections $\nabla_1$ and $\nabla_2$ to be dual, with 
\[
F_{\nabla_1} = -F_{\nabla_2},
\]
resolving orientation ambiguities.

By Theorem \ref{P:1} $\Pi_2 = \iota \circ \Pi_1$, thus the pullback bundles satisfy:
\[
E_2 = \Pi_2^* T \mathbb{CP}^2 = (\iota \circ \Pi_1)^* T \mathbb{CP}^2 = \Pi_1^* \left( \iota^* T \mathbb{CP}^2 \right).
\]
As $\iota$ is anti-holomorphic, we have:
$
\iota^* T \mathbb{CP}^2 \cong \overline{T \mathbb{CP}^2},$
so $E_2 \cong \overline{E_1}.$

\, 

The curvature $F_{\nabla_i}$ is the pullback of the curvature $\Theta$ of the bundle $\left( T \mathbb{CP}^2, g_{FS} \right)$, where $g_{FS}$ is the Fubini-Study metric. As $g_{FS}$ is Kähler-Einstein with Kähler form $\omega_{FS}$ and:
$\iota^* \omega_{FS} = -\omega_{FS},$
it follows that:
\[
\iota^* \Theta = -\Theta.
\]
Thus:
\[
F_{\nabla_2} = \Pi_2^* \Theta = (\iota \circ \Pi_1)^* \Theta = \Pi_1^* \left( \iota^* \Theta \right) = \Pi_1^* ( -\Theta ) = - F_{\nabla_1}.
\]
This antisymmetry fixes the relative signs in parallel transport, resolving orientation ambiguities.

    \item {\it Uniqueness of Solution}:
Since the projections $\Pi_i$ are immersive, the differentials:
\[
d\Pi_i : T_p O \to T_{\Pi_i(p)} \mathbb{CP}^2
\]
are injective. This implies:
\[
\ker(\omega_1) \cap \ker(\omega_2) = \{ 0 \} \subset T_p \mathbb{CP}^3.
\]
Locally, the forms $\omega_1, \omega_2$ form a full-rank system on $T_p O$. We  build a \textbf{linear system for $v\in T_p \mathbb{CP}^3$}:
\[
\begin{cases}
\nabla_1 \mu_1 = v \cdot \omega_1, \\
\nabla_2 \mu_2 = v \cdot \omega_2,
\end{cases}
\]
The above system defines a linear map:
\[
A_p : T_p \mathbb{CP}^3 \to \mathbb{C}^4, \quad A_p(v) = \left( \omega_1(v), \omega_2(v) \right).
\]
The domain of $A_p$ is $\operatorname{dom} A_p = T_p \mathbb{CP}^3$ of real dimension 6 and the codomain is of real dimension $8$. By the immersivity property, $\operatorname{rank}_\mathbb{R} A_p = 6$ (full rank) and so $\ker A_p$ is 0-dimensional. The compatibility condition $\omega_1 = \iota^* \omega_2$ ensures that $\nabla_i \mu_i \in \operatorname{im} A_p$, guaranteeing a unique solution:
\[
v_p \in T_p O \subset T_p \mathbb{CP}^3.
\]
  \item {\it Reconstruction}. The unique solution $v_p\in T_pO$ defines a smooth vector field $v$ on $O$. To reconstruct the object $O$, we introduce an initial point $p_0 \in \mathbb{CP}^3$ and solve the algebraic equations from projections and moments:
\[
\Pi_1(p_0) = z_1, \quad \Pi_2(p_0) = \iota(z_1),\quad
\mu_1(p_0) = c_1, \quad \mu_2(p_0) = c_2.
\]
These define algebraic varieties:
$V_{\Pi_i} = \{ p \mid \Pi_i(p) = z_i \} \cong \mathbb{CP}^1 \quad (\text{linear, degree 1})$,
and $V_{\mu_i} = \{ p \mid \mu_i(p) = c_i \} \quad (\text{degree } d_i \text{ in homogeneous coordinates}).$
By Bézout’s theorem, the intersection number is:
\[
[V_{\mu_1}] \cdot [V_{\mu_2}] \cdot [V_{\Pi_1}] \cdot [V_{\Pi_2}] = d_1 d_2 \times 1 \times 1.
\]
If $O$ is non-symmetric (no non-trivial automorphisms), the varieties intersect transversely, yielding a unique solution $p_0$ when $d_1 d_2 = 1$, e.g., when the $\mu_i$ are linear.

\medskip 

We discuss now the vector fields spanning $TO$.  The solution $v$ gives one vector field $v^{(1)}$. We extend it to three independent vector fields $\{v^{(1)}, v^{(2)}, v^{(3)}\}$ spanning $TO$, obtained by solving the reconstruction equations for basis vectors. We define $ v^{(2)}, v^{(3)}$ by solving {\it modified} reconstruction equations. These are obtained by rotating the moment maps (i.e. such as $\nabla_i({\bf R}_{\theta_{ik}}\mu_i)=v^{(k)}\cdot \omega_i$ where ${\bf R}_{\theta}$ is a rotation in $\mathbb{C}^2$,the indices are $i=\{1,2\}$, $k=\{2,3\}$). We choose the rotations so that $\{v^{(1)}, v^{(2)}, v^{(3)}\}$ span $T_pO$, at each $p$ and so as to ensure full rank of the linear system $A_pv^{(k)}=b^{(k)}$. Since $\operatorname{rk}_{\mathbb{R}}(A_p)=6$ and $\operatorname{dim}=3$, the system admits solutions $v^{(k)}\in T_pO$ as long as $b^{(k)}\in  \operatorname{im}(A_p)$. This holds by the curvature duality $\omega_1=\iota^*\omega_2$.

\medskip 
The vector fields $\{v^{(1)}, v^{(2)}, v^{(3)}\}$  span a 3D distribution $D\subset TO$. To reconstruct $O$ as the integral manifold of $D$, we must verify that $D$ is involutive. Due to curvature duality $F_{\nabla_1} = -F_{\nabla_2}$, we have 

\[
[v^{(i)}, v^{(j)}] \in \operatorname{span} \{ v^{(1)}, v^{(2)}, v^{(3)}\} \quad \forall i,j. 
\]

Frobenius' theorem ensures that the maximal integral manifold through $p_0\in O$ is precisely $O$ itself.

\medskip 
Finally, the reconstruction step follows from integrating flows from $p_0$ (uniquely determined by Bezout).  
\end{enumerate}
\end{proof}

\end{document}
